% Template para Trabalhos no 24°CONGRESSO BRASILEIRO DE ENGENHARIA QUÍMICA & 19° ENCONTRO BRASILEIRO SOBRE O ENSINO EM ENGENHARIA QUÍMICA
% Adaptação Realizada pela Comissão Organizadora
% Mais informações: https://www.cobeq2023.com.br/

% ----------------------------------
% IMPORTAÇÃO DE PACOTES
% ----------------------------------
\documentclass[12pt,a4paper]{article}
\usepackage[top=4.0cm, bottom=3.2cm, left=2.5cm, right=2.5cm, headheight=3.0cm,footskip=25pt]{geometry}
\usepackage[T1]{fontenc}
\usepackage{mathptmx}
\usepackage[utf8]{inputenc}
\usepackage[brazil]{babel}
\usepackage{fancyhdr}
\usepackage{ragged2e}
\usepackage{scrextend}
\usepackage{titlesec}
\usepackage{indentfirst}
\usepackage{graphicx}
\usepackage{caption}
\usepackage{textcomp} 
\usepackage{amssymb,amsfonts,amsmath}
\usepackage{enumitem}
\usepackage{xcolor}
%\usepackage[skip=6pt, indent=1cm]{parskip}
\usepackage[alf,abnt-etal-list=2]{abntex2cite}
% IMPORTE PACOTES ADICIONAIS, CASO PRECISE
%\usepackage{....}

% ------------------------------------
% DEFINIÇÕES DO TEMPLATE - NÃO ALTERAR
% ------------------------------------
\makeatletter
% Afiliação
\def\@affiliation{}
\newcommand{\affiliation}[1]{
  \def\@affiliation{#1}
   }
% Resumo
\def\@abstract{}
\renewcommand{\abstract}[1]{
  \def\@abstract{#1}
   }
% Seção
\titleformat{\section}
  {\normalfont\fontsize{14}{15}\bfseries}{\uppercase{\thesection}}{1em}{}
  
\titlespacing*{\section}{0pt}{6pt}{0pt}
\titlespacing*{\subsection}{0pt}{6pt}{0pt}

% Títulos, autores e resuno
\renewcommand\maketitle{
   \vspace*{-1cm}
   {\noindent\Large\bfseries\@title\justifying}\\
   \begin{center}
   {\@author}   
   \end{center}
   
   \begin{center}
   {\setlength{\parskip}{0pt}
   \@affiliation
   }
   \end{center}

    \vspace*{\baselineskip}
   
    \begin{addmargin}[1cm]{1cm}
    {\itshape\noindent RESUMO-\@abstract}
    \end{addmargin}
    
  }
\makeatother

\setlength{\parindent}{1cm}
\setlength{\parskip}{4pt}

% LEGENDAS
\captionsetup{justification=centering, textfont=it, labelfont=it, skip=4pt, aboveskip=4pt}

% ESTILO DA BIBLIOGRAFIA
\bibliographystyle{abntex2-alf}

% CABEÇALHO E RODAPÉ
\definecolor{greentemplate}{RGB}{48, 154, 90}
\renewcommand{\footrulewidth}{1.5mm}
\let\FootRule\footrule
\renewcommand\footrule{\textcolor{greentemplate}{\FootRule}}
\pagestyle{fancy}
\setlength{\headheight}{3.92cm}
\lhead{\includegraphics[width=\textwidth]{cabecalho-2023}}
\chead{}
\rhead{}
\cfoot{}
\rfoot{\thepage}
\renewcommand{\headrulewidth}{0pt}

% ----------------------------------
% DADOS DO TRABALHO
% ----------------------------------
%Título
\title{O TÍTULO DO TRABALHO COMPLETO DO XXIII CONGRESSO BRASILEIRO DE ENGENHARIA QUÍMICA DEVERÁ SER EM FONTE TIMES NEW ROMAN 16, NEGRITO, JUSTIFICADO, ESPAÇO SIMPLES}
% Autores
\author{AB SOBRENOME$^1$, C de SOBRENOME$^2$, D SOBRENOME Jr$^3$, E SOBRENOME$^4$ (Times New Roman 12)}
% Afiliação dos autores
\affiliation{%
$^1$ Universidade Federal da Bahia, Escola Politécnica, Departamento de Engenharia Química

$^2$ SENAI CIMATEC — Centro Integrado de Manufatura e Tecnologia

$^3$ Universidade Salvador

$^4$ Instituto Federal de Educação, Ciência e Tecnologia - Bahia

E-mail para contato: nome@wxyz.br
}
%Resumo
\abstract{A palavra RESUMO deverá ser digitada em letras maiúsculas e o texto em letras minúsculas, em itálico, continuando na mesma linha em um único parágrafo com letra Times New Roman 12 e espaçamento simples. Deverá haver um espaço de 1 cm depois da margem esquerda e 10 mm antes da margem direita. O texto deverá conter o objetivo do trabalho, a metodologia empregada e os principais resultados quantitativos e qualitativos. O texto não deverá conter referências bibliográficas.}

% -------------------------------
% INÍCIO DAS SEÇÕES DO DOCUMENTO
% -------------------------------
\begin{document}

\maketitle

\section{INTRODUÇÃO}

O título da seção é um cabeçalho de primeira ordem, deverá ser numerado com algarismos arábicos, alinhados junto à margem esquerda, com letras maiúsculas e em negrito. Deverá ser separado por um espaçamento de 6 pt antes e depois do texto. As seções incluídas neste template devem ser utilizadas como orientações para escrita do trabalho, devendo os autores definirem a quantidade de seções e seus respectivos títulos.

O texto deverá ser digitado em espaço simples para todo o trabalho e espaço de 4 pt entre parágrafos, contendo entre 4 e 6 páginas, com até 8 autores. A página deverá ser configurada no formato A4 (210mm x 297mm), com orientação ``retrato'', em uma coluna, com margem superior 40 mm e margens inferior, direita e esquerda de 25 mm. A fonte deverá ser a Times New Roman, tamanho 16 no título, 14 nos cabeçalhos de primeira e segunda ordem, 12 no texto e cabeçalhos de terceira ordem e 12 nas notas de referência dos autores abaixo do título. Recomenda-se que o trabalho completo seja escrito neste mesmo arquivo fornecido pela comissão organizadora.

Os trabalhos redigidos em português, inglês ou espanhol deverão ter a mesma formatação. A escrita dos nomes e símbolos das unidades, e a expressão dos valores das grandezas deve seguir o Sistema Internacional de Unidades, SI, por exemplo, em trabalhos em português o separador decimal é a vírgula, para trabalhos em inglês é o ponto.  A primeira linha de cada parágrafo deverá iniciar com uma tabulação de 1 cm da margem esquerda, com alinhamento justificado. 

\section{METODOLOGIA}

\subsection{Segundo Cabeçalho}

Os cabeçalhos de segunda ordem deverão ser alinhados junto à margem esquerda, digitados em negrito, sendo maiúscula apenas a primeira letra de cada palavra com separação de um espaço do texto anterior.

\textbf{\underline{Terceiro cabeçalho:}} O cabeçalho de terceira ordem inicia após tabulação de 1 cm da margem esquerda, sem numeração, com caracteres sublinhados, sendo maiúscula apenas a primeira letra da primeira palavra. O texto deverá começar na mesma linha do cabeçalho.

\section{RESULTADOS E DISCUSSÃO}

\subsection{Figuras}

As figuras poderão ser coloridas e deverão ser inseridas no corpo do trabalho, tão próximas quanto possível das citações, no texto, sobre elas. As figuras deverão ser centralizadas, sem exceder o tamanho limitado pelas margens da página.

Cada figura deverá ter um título numerado em algarismos arábicos. Os títulos deverão ser centralizados na parte superior da figura, separados por espaço de 4 pt antes e após o texto e digitados como: Figura 1 – Título da figura com ponto final. No texto, elas deverão ser mencionadas, por exemplo, da seguinte forma: ``conforme mostra a Figura \ref{fig:figura-exemplo}...''.

\begin{figure}[!hb]
    \centering
    \caption{Isoterma de adsorção do componente xxx ($\bullet$)  e do componente yyyy ($\blacksquare$)  em partículas porosas de zzzz a temperatura de 25 \textdegree C.}\label{fig:figura-exemplo}
    \includegraphics[width=0.6\textwidth]{Figura-exemplo}
\end{figure}

\subsection{Tabelas}

As tabelas deverão ser centralizadas e posicionadas após e próximas às suas respectivas citações no texto, com os títulos centralizados na sua parte superior com espaço de 4 pt antes e após o texto e digitados como: Tabela 1 – Título da tabela sem ponto final. As tabelas deverão ser mencionadas no texto, por exemplo, da seguinte forma “...o resultado do planejamento pode ser encontrado na Tabela 1”. As unidades de medida correspondentes a todos os termos deverão ser claramente indicadas, atentando-se ao S.I.

\begin{table}[htb]
    \centering
    \caption{Propriedades texturais dos materiais A, B e C.}
    \label{tab:tabela-exemplo}
    \begin{tabular}{cccc}\hline
         Materiais & $S_{\text{BET}}$/(m\textsuperscript{2}/g) & $V_{p}$ / (cm\textsuperscript{3}/g) & $D_{p}$ / nm  \\\hline
         A & 80,2 & 0,45 & 24,0\\
         B & 65,3 & 0,48 & 24,6 \\
         C & 38,3 & 0,56 & 17,0\\\hline
    \end{tabular}
\end{table}

\subsection{Equações}

As equações deverão ser escritas em itálico, centralizadas, com numeração consecutiva entre parênteses, rente à margem direita, conforme exemplo:
\begin{align}
    F = K\cdot x. \label{eq:equacao-exemplo}
\end{align}

Deverão ser mencionadas no texto, por exemplo, da seguinte forma: ``... a substituição da Equação \eqref{eq:equacao-exemplo} ...''. Equações com mais de uma linha deverão ser numeradas na última linha, entre parênteses e rente à margem direita.

\section{CONCLUSÃO}

As principais conclusões a respeito do trabalho devem ser apresentadas de forma sucinta e clara.

\section{NOMENCLATURA}

Caso necessário, a lista de notações e símbolos utilizados, assim como suas unidades de medida, seguindo-se o Sistema Internacional - SI, deverão ser relacionados antes das referências bibliográficas por ordem alfabética.

\section{REFERÊNCIAS}

As referências bibliográficas deverão ser citadas no texto, através do último sobrenome do autor e do ano de publicação, o qual deverá estar entre parênteses. Neste template usamos o pacote abntex2cite para gerir as referências, conforme os exemplos a seguir: “Os trabalhos de \citeonline{Ferreira2016}, \citeonline{Costa2000} e \citeonline{Vasquez2000} mostraram...”, ou “... tem sido mostrado \cite{Vasquez2000,Costa2000,Ferreira2016}”. No caso de dois autores, ambos deverão ser citados. O pacote abntex2cite realizada este controle, exemplo: “... segundo \citeonline{Vasquez2000}...” Em caso de três ou mais autores, deverá ser citado o sobrenome do primeiro autor seguido da expressão “et al.”.

A lista de referências deverá incluir somente os trabalhos citados no texto, relacionados em ordem alfabética, de acordo com o sobrenome do primeiro autor, iniciando na margem à esquerda, com tabulação de 1 cm a partir da segunda linha e separadas por espaço simples. As referências deverão conter todos os sobrenomes e iniciais dos autores, separados por vírgula. Os títulos dos periódicos deverão aparecer abreviados e em itálico, conforme o exemplo: %
% Mantenha este código
\begingroup
\renewcommand{\section}[2]{}%
% caso precise, altere o nome do arquivo que contém as referências
\bibliography{Referencias}
\endgroup

\end{document}
